% Part 2 subsection (human-written voice; Andy edits): what the autonomous loop exploited, and how it was caught.
% Numbers marked [TBD] are filled from results/chimera_audit_gpu.json, runs/hyak/eval_refix*.json and run 9.
\subsection*{2.x The loop exploited a defect in its own harness}
\textbf{What happened.} The harness keeps $K$ archive slots that testers save to and restore from. When several of the $B$
vectorised workers chose the same slot in one step --- which the discovered program's ``replace the pool's worst'' speed and
proximity pools do whenever many workers beat the incumbent --- the state, its feature vector and its lineage were written by
separate scatter operations, and on the GPU those scatters did not agree on a winner. The slot could then hold a
\emph{chimera}: the positions of one world with the velocities of another. Restoring such a state and stepping it produces
violent, physically impossible motion, and OOB, ANOMALY and TUNNEL witnesses fire on exactly that. The program search
saw only the fitness, and it selected for programs that restore from those pools as often as possible.

\textbf{How it was caught.} Not by the fitness, which rewarded it, and not by the transfer test, which it passed: chimera
restores work on any level and any class. It was caught by the requirement that every real-anomaly report be reproducible
from a level reset by replaying its action chain. Chains reconstructed from the new append-only lineage log were complete
yet replayed only 12 of 86 flagged runs, on the GPU and on the CPU alike. A lineage audit that replays the chain of every
archived state and compares it with the stored state reproduced the defect in seconds (13 of 64 archived states
mismatched, all after restores) and confirmed the fix (0 of 64, 0 of 96, 0 of 64). The fix keeps exactly one writer per
slot per step.

\textbf{What changed.} On the fixed harness the discovered program scores 0.568 / 60\% on the search split (three seeds) against action-repeat
random's 0.565 / 59\%, not 0.666 / 72\%; on all three splits with fresh seeds (100--104) the discovered program scores 0.574 / 0.557 / 0.267 against action-repeat's
0.569 / 0.579 / 0.266 and Go-Explore's 0.537 / 0.556 / 0.220 --- level with action-repeat everywhere; the searches re-run on the
fixed harness reach [TBD run 9/10]. The GPU audit measures the mechanism directly on 16 search instances at $128\times2048$: under the
legacy race the program finds 62\% of targets (ANOMALY 75\%) and 24\% and 7\% of its archived states on the two audited
instances are chimeras --- their stored feature vectors disagree with the features of the stored state; under the fixed harness
it finds 50\% (ANOMALY 50\%) and no archived state is a chimera. On the CPU the same duplicate-index scatters resolve
consistently and the audit shows no chimeras and no score difference, which is why local checks never caught it.

\textbf{What we take from it.} (i) A discovery loop that scores itself will find the cheapest path to a high score, and a
defect in the evaluator is the cheapest path there is; the loop's own diagnostics (transfer, ablations, portfolios) all
confirmed the exploit because they all ran on the same evaluator. (ii) The check that caught it was the one that did not
go through the evaluator: independent reproduction of a concrete artefact. Every claim in this paper that survives is one
that has such a check. (iii) The defect is ordinary vectorised-JAX code --- duplicate scatter indices --- and we found no
mention of it in the exploration-testing literature we build on; testers that archive and restore states in parallel
should audit their archives the way we now do (Part~2, \texttt{scripts/lineage\_audit.py}).
